\documentclass[12pt]{article}
\usepackage[margin=1in]{geometry}
\usepackage[T1]{fontenc}
\usepackage{xcolor}
\definecolor{garnet}{HTML}{73000A}

\setlength{\parindent}{0pt}
\setlength{\parskip}{0.5em}

\pagestyle{empty}
\begin{document}

\noindent
Department of Mechanical Engineering\\
University of South Carolina\\
Columbia, SC 29208\\
\today

\noindent
{\color{garnet}\rule{\textwidth}{0.5pt}}

\vspace{0.3em}

\noindent
Dear Members of the Evaluation Committee,

I enthusiastically recommend Noah Young for the NASA SC Space Grant Undergraduate Research Award. As his faculty advisor in the Integrated Multiphysics and Systems Engineering Laboratory (iMSEL) since August 2025, I have supervised his research on radar and autonomous perception systems from his first few weeks of college.

Measured against the broader pool of students I have taught, Noah ranks in the top five percent. He joined iMSEL at the start of his first semester as an Honors College aerospace engineering student and is now contributing to a planned ASME IDETC 2026 submission on synthetic radar data augmentation, in which he developed a method using land/water segmentation and object extraction to regenerate radar imagery from real sensor data.

When he began, radar fundamentals were almost entirely new to him. Rather than waiting for instruction, he independently read background literature on radar captures, pulse compression, and geometric simplification methods, then applied what he learned to produce working code within weeks. Somewhat impressively, he has maintained a 4.0\,GPA while carrying this research workload.

His technical contributions span multiple systems in our lab. As a first task, he fine-tuned a YOLO model for maritime small-object detection and deployed it on the NVIDIA Jetson Orin -- beginning some of our earlier tests on real-time radar object detection. He additionally designed a 3D-printable radar antenna mount in CAD that enables our Furuno NXT radar to be integrated onto a research vessel for offshore data collection. He also built a raw radar-to-PNG conversion pipeline and real-time visualization tools, in Rust, that the radar team now uses for post-processing and computer vision workflows. In each case he went beyond the assigned task, identifying and solving adjacent problems without being asked.

His proposed project -- an open-source terrain-occluded radar simulator validated against measured returns -- grows directly from this work. He has demonstrated he can learn an unfamiliar domain independently and deliver functioning software on schedule, and I believe the simulator is scoped appropriately for him.

I wholeheartedly recommend Noah Young for this award.

\vspace{0.5em}

\noindent
Sincerely,

\vspace{1em}

\noindent
Dr.\ Yi Wang\\
Associate Professor\\
Department of Mechanical Engineering\\
University of South Carolina\\
\textit{yiwang@cec.sc.edu}

\end{document}
